\section{机器翻译评估}

本讲的第一部分承接上一讲的序列到序列（Sequence-to-Sequence）机器翻译模型，讨论如何评估机器翻译系统的质量。

\begin{figure}[H]
\centering
\includegraphics[width=0.95\textwidth]{slides-images/slide-001.jpg}
\caption{本讲课程大纲：涵盖 MT 评估、注意力机制、期末项目指导及实践建议\protect\footnotemark}
\end{figure}
\footnotetext{Slides 第2页。}

\subsection{序列到序列模型回顾}

在上一讲中，我们介绍了基于多层 LSTM 的编码器-解码器（Encoder-Decoder）架构用于机器翻译。编码器读入源语言句子，将信息压缩到最后一个隐藏状态中；解码器以此为初始状态，逐词生成目标语言翻译。

\begin{figure}[H]
\centering
\includegraphics[width=0.95\textwidth]{slides-images/slide-002.jpg}
\caption{多层深度编码器-解码器机器翻译网络：源句（德语）通过多层 LSTM 编码，解码器使用不同参数逐词生成翻译（英语）。底部标注``Conditioning = Bottleneck''指出了该架构的核心问题\protect\footnotemark}
\end{figure}
\footnotetext{Slides 第3页。参考 Sutskever et al. 2014; Luong et al. 2015。}

\begin{warningbox}{信息瓶颈问题}
在标准的 Seq2Seq 架构中，源句的\textbf{所有信息}必须被压缩到编码器最后一个隐藏状态向量中。对于短句子这或许可以接受，但当源句长达 40 个词时，将全部语义信息塞入一个固定大小的向量显然是不合理的。这个问题被称为\textbf{信息瓶颈（Information Bottleneck）}，也是注意力机制被发明的直接动因。
\end{warningbox}

\subsection{BLEU 评估指标}

为了自动评估机器翻译质量，Papineni 等人（2002）提出了 \textbf{BLEU}（\textbf{B}ilingual \textbf{E}valuation \textbf{U}nderstudy）指标。在此之前，翻译质量只能通过人工评估来衡量，这虽然是金标准，但成本高昂且无法嵌入训练循环。

\begin{figure}[H]
\centering
\includegraphics[width=0.95\textwidth]{slides-images/slide-003.jpg}
\caption{BLEU 评估方法：将机器翻译与一个或多个参考翻译进行 n-gram 匹配\protect\footnotemark}
\end{figure}
\footnotetext{Slides 第4页。来源：Papineni et al., 2002。}

\begin{importantbox}{BLEU 指标的核心思想}
BLEU 通过计算机器翻译与人工参考翻译之间的 \textbf{n-gram 重叠度}来衡量翻译质量：
\begin{itemize}
    \item 计算 1-gram、2-gram、3-gram、4-gram 的精确率（precision）
    \item 取这些精确率的\textbf{几何平均值}
    \item 加入\textbf{简短惩罚}（brevity penalty），防止系统通过只翻译简单部分来获取高分
\end{itemize}
BLEU 分数理论上在 0--100 之间，但由于翻译的多样性，永远不会达到 100。
\end{importantbox}

\begin{figure}[H]
\centering
\includegraphics[width=0.95\textwidth]{slides-images/slide-004.jpg}
\caption{BLEU 评分示例：机器翻译与4个参考翻译的 n-gram 匹配过程，圈出的词组表示匹配成功的片段\protect\footnotemark}
\end{figure}
\footnotetext{Slides 第5页。}

\begin{knowledgebox}{BLEU 分数的解读}
\begin{itemize}
    \item \textbf{20 分左右}：可以大致理解源文档的主题
    \item \textbf{30--40 分}：翻译质量开始变得不错
    \item \textbf{50--60 分}：现代神经机器翻译系统在一些语言对上已能达到的水平
\end{itemize}
原始设计建议使用多个参考翻译来更好地覆盖翻译空间，但实践中使用单个参考翻译也很常见。
\end{knowledgebox}

\begin{warningbox}{BLEU 的局限性}
BLEU 是一个\textbf{有用但不完美}的指标：
\begin{itemize}
    \item 一个好的翻译可能因为措辞与参考翻译不同而获得低分
    \item 一个差的翻译可能因为碰巧包含一些正确的词而获得一些分数
    \item 它无法捕捉语义等价性——同义词替换会被错误惩罚
\end{itemize}
因此，人工评估在高风险场景中仍然是必不可少的。
\end{warningbox}

\subsection{机器翻译的历史进展}

\begin{figure}[H]
\centering
\includegraphics[width=0.95\textwidth]{slides-images/slide-005.jpg}
\caption{机器翻译 BLEU 分数随时间的变化（Edinburgh En-De WMT）：统计短语翻译（红色）进展缓慢，句法翻译（紫色）略有提升，神经机器翻译（青色）从2015年起呈现爆发式增长\protect\footnotemark}
\end{figure}
\footnotetext{Slides 第6页。}

机器翻译技术经历了三个重要阶段：

\begin{enumerate}
    \item \textbf{统计短语翻译}（Phrase-based SMT，1990s--2010s）：IBM 开创，Google Translate 最初使用此方法。进展缓慢，BLEU 分数几乎停滞。
    \item \textbf{句法翻译}（Syntax-based SMT，2005--2015）：试图利用句法结构改善翻译，尤其针对词序差异大的语言对（如德-英、中-英），但改进幅度有限。
    \item \textbf{神经机器翻译}（Neural MT，2014至今）：2014年首次出现，2015年在 bake-off 评测中亮相，2016年超越所有传统方法，此后一路高歌猛进。
\end{enumerate}

\begin{knowledgebox}{NMT 的崛起}
Manning 教授指出，在 2005--2014 年间，机器翻译社区的主流观点是``要想做好翻译，必须理解句法结构''。然而，神经机器翻译的成功表明，端到端学习可以隐式地捕获这些结构信息，这彻底颠覆了传统思维。注意力机制是推动 NMT 成功的关键技术之一。
\end{knowledgebox}

\subsection{本章小结}

\begin{itemize}
    \item BLEU 是最广泛使用的机器翻译自动评估指标，基于 n-gram 精确率的几何平均
    \item BLEU 有用但不完美，好翻译可能低分，差翻译可能得分
    \item 神经机器翻译从2014年起迅速超越传统方法，注意力机制是其成功的核心要素
    \item 编码器-解码器架构存在信息瓶颈，这直接催生了注意力机制的发明
\end{itemize}

%% ========== Part 2: 注意力机制 ==========

